\chapter{Context: languages and tools}

At the very base of my PhD there are two softwares, the Rocq Proof assistant
(formerly known as \coq~\cite{coqart}) and the \elpi~\cite{dunchev15}.

\section{The \rocq proof assistant}

\cite{coqart}

Theorem provers with HO logic programming languages:~\cite{felty88}

Mathematical components: \cite{assia2022}

Abella: \cite{abella}

SSReflect: \cite{ssreflect}

Isabelle: \cite{isabelle}

\subsection{The basis of the language}

Function vs Propositions

\subsection{Elaboration}

As explained in~\cite{elaboration}, raw text is the mean bridging the communication
between the and the theorem prover. This 

\begin{figure}
  \includegraphics[width=1\textwidth]{elaboration.pdf}
  \caption{Elaboration skeme}
  \label{fig:elaboration}
\end{figure}

In \cref{fig:elaboration} (taken from~\cite{elaboration}), we have a skematical
representation of the elaboration process. \cite{elaboration}

\paragraph{Typechecking} The rich type theory of the Coq system confers to type
inference the power of computing values and proofs. 

\paragraph{Coercions} \cite{barthe95}: Used in general to forget information
about an object



\subsection{Ad-hoc polymorphism}

\paragraph{A bit of history}

Polymorphism is a fundamental feature of functional programming languages. In
his lecture notes, Strachey~\cite{polymorphism} distinguishes between two main
forms of polymorphism: parametric polymorphism and ad-hoc polymorphism.
Parametric polymorphic functions are \emph{uniform}, meaning that their behavior
does not depend on the specific type at which they are instantiated. For
example, the function \emph{length}, with type \mlI{'a list -> int}, returns $7$
when applied to any list of seven elements. In other words, the compiler
reuses the same implementation of \emph{length} for every call, regardless of
the instantiation of \mlI{'a}.

Parametric polymorphism is widely used and highly practical in functional
programming. However, it is not suitable for modeling mathematical
function overloading. In mathematics, the symbol $+$ has a consistent,
intuitive meaning, and no ambiguity arises in expressions such as $1 + 2$
or $\pig + e$. In the first expression, $+$ denotes addition over integers;
in the second, it denotes addition over real (or rational) numbers. The same
symbol is consistently used for operations that share common algebraic
properties.

In contrast, the simply typed lambda calculus does not allow a single symbol to
have multiple type-dependent interpretations. This limitation appears in
practical languages such as \ocaml, where integer and floating-point additions
are represented by distinct operators, $+$ and $+.$, respectively. Thus, one
writes $1 + 2$ for integers and $\pig +\!.\ e$ for floating-point numbers. In
general, a different symbol must be introduced for each type implementing a
given mathematical operation.

To address the need for a shared notation for operations that sharing common
properties across types, Wadler and Blott introduced the notion of type classes
in 1989~\cite{wadler1989-tc}. Type classes provide a principled form of ad-hoc
polymorphism, enabling a single symbol to be associated with multiple
type-specific implementations while preserving static type safety.

\paragraph{Type classes}
A type class is a structure parametrized by one (or more) types $t$ and contains
the signatures of a set of functions (called \emph{methods}) that may depend on
$t$. A type class serves as an interface that its instances are required to
implement.


They are implemented in Haskell for the first time and then in coq, in coq
they also bring proof and propositions

Priorities

Type-class resolution is part of the elaborator and, since it is an key aspect
of this PhD, we dedicate a full section to it.

The notion of type class has been introduced in~\cite{wadler1989-tc} has a way
to overload arithmetic operators. Differently from \emph{parametric}
polymorphism, where a 


Used to enrich the information about an object.

Programming in a way reminiscent of Prolog, used by Gonthier in Odd Order Theorem.

Extend the unification alogirhm of the coq system, looking in a databae of hints

In the elaboration phase, one of the 

Wadler: \cite{wadler1989-tc}

Matthieu: \cite{sozeau2008-tc}

Haskell: \cite{wenzel97-tc}

ATTENZIONE: Eqb is not eqb

We can define a type class for testing equality. In \cref{fig:eqb-tc}, we give
two implementations of \texttt{Eqb} in the Rocq and Haskell languages. The
similarity between the two pieces of code is striking, but in Rocq we can go
further. Type classes can be used not only for function overloading, but also
for theorem overloading. In coq a class can in fact not also carry function
specifications, but also the porperties that these functions should satisfy. For
example, we could have a method, say \emph{eqP} certifying that for
each implementation of \emph{eqb} there is a proof that the two objects are
structurally equivalent. This would also mean that each implementation of
\emph{Eqb} should implement \emph{eqP}.

This last implementation with eqb and eqP is the one used by
ssreflect~\cite{ssreflect} through~\cite{hb}

Type classes are deppely used in libraries such as stdpp, iris, ...

Cite Haskell~\cite{wadler1996-hs}, in caml it does not exists: \texttt{+.} vs
\texttt{+}

IRIS:~\cite{iris-junk}

\begin{figure}
  \begin{subfigure}{1\textwidth}
  \includegraphics[width=1\textwidth]{./code/base_tc.pdf}
  \caption{\texttt{Eqb} typeclass in Rocq}
  \label{fig:eqb-tc-coq}
  \end{subfigure}
  \begin{subfigure}{1\textwidth}
  \includegraphics[width=0.75\textwidth]{./code/base_tc_hs.pdf}
  \caption{\texttt{Eqb} typeclass in Haskell}
  \label{fig:eqb-tc-hs}
  \end{subfigure}
  \caption{Eqb typeclass in different languages}
  \label{fig:eqb-tc}
\end{figure}



\paragraph{Canonical structures}: \cite{cs}
Used to enrich the information about an object. In mathComp


% \begin{minted}{coq}
% AA
% \end{minted}


\paragraph{Typeclasses vs Canonical structures}

``Find me a structure that works'' $\to$ typeclasses\\
``This type already carries its structure'' $\to$ canonical structures

\section{The \elpi language}

\subsection{Operational semantics}

\paragraph{Backchaining}
  % $$
  % \!\! \mathcal{F}(\prog, \pred\, \vecL{t}, \vecL{g}, \subst, a) =
  % \bigg[
  %   %\bigg(
  %   \subst['],
  %   %\Big(
  %   %\underbrace{
  %     \big(\big[(\prog, g, a) \mid g \in\! \vecL{b}\big] \atsign\ \vecL{g}\big)
  %   %}_{\mathrm{premises}} \atsign ~ \vecL{g}
  %   \ \kIF\ \mathcal{H}(\vecL{tu}, \subst) = \subst['] \not= \bot
  %   %\Big)
  %   %\bigg)
  %   ~\bigg\rvert~
  %   (\clauseCmd{\pred}{\vecL{u}}{\vecL{b}}) \in \prog\ \pred
  %   \bigg]
  % $$
  % $$
  % \mathcal{H}(\vecL{tu}, \subst) = \fold\ \unify\ \vecL{tu_o}\ (\fold\ \match\ \vecL{tu_i}\ \subst)
  % $$

\begin{figure}[h]
  \centering
  \includegraphics{img/det_check/dynamic.pdf}
  \caption{Elpi semantics}
\end{figure}

%

\paragraph{Controlling backtracking: the \texorpdfstring{$!$}{cut} operator}

\paragraph{Higher-order programming}

Elpi~\cite{dunchev15} is a dialect of $\lambda$prolog~\cite{miller1991-pf,miller1988}

Teyjus:~\cite{teyjus,miller2012}

Elf: \cite{elf}

Twelf:\cite{twelf}

Lambda prolog: \cite{qi2009}

Mercury: \cite{somogy1996}

WAM: \cite{ait1991-wam}

Prolog semantics: \cite{debray1988}

HOAS: \cite{pfenning1988}

\subsection{Higer-order unification}


\paragraph{Occur check}

\paragraph{Unification in the pattern-fragment (\llambda)}

Miller: \cite{miller1991-pf}


Huet: \cite{huet75}

Modes: \cite{zhou2012}


SWI prolog: \cite{swi-prolog-cut}

Andrews:~\cite{andrews2003}

Vink: \cite{vink1989}

\subsection{\chr: goal suspension}

\cite{chr,chr15}

CHR elpi: \cite[section 4.3]{chr-elpi}

CHR Pfenning: \cite{pfenning1993}

% \section{\coqelpi: \rocq and \elpi interleaving}
\section{\coqelpi: a logical framework for \rocq}

% \subsection{Object language (\Fo{}) and Meta programming (\Ho{})}
Distinction between object and meta-language

\subsection{HOAS: \texorpdfstring{$\lambda$}{lambda}-terms representation}

In order to implement such a search in Elpi, we shall describe the encoding
of Coq terms and then the encoding of instances as rules.
Elpi comes equipped with
a Higher-Order Abstract Syntax (HOAS~\cite{pfenning1988}) datatype of
Coq terms, called \eI{tm}, that includes (among others) the following
constructors:

\begin{elpicode}
type lam  tm -> (tm -> tm) -> tm.     % lambda abstraction
type app  list tm -> tm.              % n-ary application
type all  tm -> (tm -> tm) -> tm.     % forall quantifier
type con  string -> tm.               % constants
\end{elpicode}

\subsection{Unification in \Fo{} and \Ho{}}

In order to reason about unification we provide a description of the
\Fo{} and \Ho languages where unification variables
are first-class terms, i.e. they have a concrete syntax as
shown in \cref{code:common-terms}.
Unification variables
in \Fo{} (\eI{o-uva} term constructor) have no explicit scope:
the ``arguments'' of a higher-order variable are given via the \eI{o-app}
constructor. For example the term <<\cI{P x}>> is represented as
<<\eI{o-app [o-uva N, x]}>>, where \eI{N} is the memory address
of \eI{P} and \eI{x} is a bound variable.
In \Ho{} the representation of <<\cI{P x}>> is instead <<\eI{uva N [x]}>>,
since unification variables are higher-order and come equipped with an
explicit scope.

Coqelpi: \cite{tassi2018}

Hierarchy Builder: \cite{hb}

Trackt: \cite{enzo2023}

Equality test~\cite{tassi2023}

Derive: \cite{tassi2019}


\subsection{Concluding remarks}

Overall, Elpi is good for elaboration
